\newpage
\begin{center}
  \textbf{\large ЗАКЛЮЧЕНИЕ}
\end{center}
\refstepcounter{chapter}
\addcontentsline{toc}{chapter}{ЗАКЛЮЧЕНИЕ}

В рамках данной дипломной работы была разработана и апробирована универсальная методика анализа защищённости мобильных приложений с использованием прокси-сервера для перехвата и анализа сетевых вызовов. Основные результаты работы заключаются в следующем:
\begin{enumerate}

\item Разработана методика анализа сетевого трафика мобильных приложений, позволяющая выявлять уязвимости, связанные с недостаточной защитой соединений, слабой аутентификацией и утечкой конфиденциальной информации. Методика успешно протестирована на трёх различных приложениях: для заказа еды, доступа к бизнес-клубу и игры Minecraft, что подтвердило её универсальность и применимость к различным типам приложений и архитектурам.

\item Установлены ключевые недостатки современных мобильных приложений в области безопасного обмена данными, включая отсутствие асимметричного шифрования и проверки сертификатов, что создаёт риски атак типа «человек посередине». На основе анализа предложены рекомендации по повышению защищённости, включая внедрение сквозного шифрования, закрепление сертификатов, минимизацию передаваемых данных и использование современных протоколов аутентификации.

\item На основе собранных данных разработан отечественный аналог приложения для заказа еды с открытым исходным кодом. Реализация альтернативного клиента и замена уязвимых модулей продемонстрировали возможность создания безопасных и конкурентоспособных решений, соответствующих требованиям локального рынка. Разработанный прокси-сервер для постоянного мониторинга сетевого трафика позволил автоматизировать анализ безопасности и выявлять потенциальные угрозы в реальном времени.


\end{enumerate}

Разработанная методика может быть использована для дальнейшего анализа защищённости мобильных приложений, создания безопасных программных продуктов и повышения уровня кибербезопасности в условиях цифровизации и импортозамещения. Работа вносит вклад в развитие отечественных технологий в области мобильной разработки и обеспечения информационной безопасности.
